\documentclass[11pt]{article}
\usepackage[a4paper,margin=24mm]{geometry}
\usepackage{amsmath,amssymb,amsthm,booktabs}
\usepackage[colorlinks=true,urlcolor=blue]{hyperref}
\newtheorem{proposition}{Proposition}
\newcommand{\HH}{\mathcal H}
\newcommand{\hh}{\mathcal h}
\newcommand{\alow}{\mathfrak a}
\newcommand{\Lam}{\Lambda}
\title{Insertion of the certified finite table into the eleven-term functional}
\author{Independent local mathematical audit}
\date{8 October 2026}
\begin{document}
\maketitle

This note proves state admission and the resulting integral prescriptions,
conditional on the primitive signed-prefix distribution, positive-product
switching and finite-seed inputs explicitly retained in
\texttt{finite-step-typed-closure-audit.tex}. It does not certify those
primitive analytic theorems, a numerical value of the whole functional,
or a new Goldbach exponent. It uses the independently replayed 200 closed
cells on $[1,5]$, 24 generations, $K_0=681/100$, and initial seed
$2327/100000$ in \texttt{certify\_finite\_step\_table.json}.
It does not use the broader floating $M=10$ table or a new second-weighted
parameter bank. Source hashes and exact checks accompany this note.

\section*{State and cell conventions}
Write $A(s)=sF(s)/(2e^\gamma)$ and $\alow(s)=sf(s)/(2e^\gamma)$;
$A=1$ on $[1,3]$ and $\alow=0$ below two. Let $q_i=1+i/50$,
$0\le i\le200$. The saved integers divided by $2^{80}$ are
uniform bounds $H_i,h_i$ throughout $[q_i,q_{i+1}]$.
Define $\HH,\hh$ by these constants in cell interiors; at a shared
endpoint take the minimum of the two constants. Put them equal to zero
outside $[1,5]$, and put $\hh=0$ below two. Interpolation is not allowed.
A parameter box meeting several cells uses the minimum of their bounds;
a box meeting the complement of the certified support can use zero.

An admitted state retains the actual marked prefix $B=N^T$, one fixed
total level $Q=N^\nu$, residual exponent $L=\nu-T>0$, ratio
$\alpha=(1-T)/L\ge2$, branch labels, and a physical ceiling $\chi$.
For an enhanced branch every original full square test and every
selected-remaining square test must be strictly below $\nu$, while
$\nu$ is strictly below the branch cap and above its gate.
The largest specified slots and any selected subset remain marked.
For a future-third-slot $W$ bound $c$, also require $\chi\le c$.
The actual physical cutoff $L/s$ must be at most $\chi$.
Classical states instead use a fixed $\nu<1/2$ and complete-modulus
support; no enhanced full-square restriction is imported into them.

The previously audited finite actions preserve these data: an appended
prime of exponent $x\le\chi$ with $2x<L$ has residual $L-x$ and
$\alpha'=1+(\alpha-1)/(1-x/L)\ge\alpha$. Each new full and
selected-remaining test is $T+2x<\nu$; subsequent tests come from the
residual well support. Old labels, signs and multiplicities stay intact.
In particular a lower child need not admit the initial seed directly
at its own argument: its actual cutoff, and the cutoffs after the
licensed descent, are what must be admitted. These facts are used here
only for the certified $M=5$ actions, not for arbitrary numeric tables.

For the main statement assume the original sufficient threshold range
\[
 1/12<a<h<d<e=57/224,\quad w_*=25/224<h<\min(\mu,2a),
 \quad 1/3<h+d<1/2,\quad 4a+d<1,\quad4h+d<1,
\]
where $\mu=25/128$, $\Lam=153/256$, $\nu_*=249/448$ and
$v(t)=\min((1+t)/2,57/64-3t/2)$.
The separate positive-$a$ extension can also be used with its primitive
inputs, but requires the extra ratio and table-support guards below.

\section*{The one-prefix terms}
The ordinary restricted lower term is retained exactly:
\[
 L_1=\frac4{\nu_*}\alow(\nu_*/a).
\]
It is not a typed-table state just because it has a sieve argument.
With empty marked prefix at its enhanced level, its ratio would be
$1/\nu_*=448/249<2$. Also its restricted atoms have varying largest
rough slots rather than the fixed marked data of the finite-table
theorem. Thus no direct $\hh(\nu_*/a)$ substitution is authorized.

For $L_2$, the exact least-prime identity starts at $w_*$ and subtracts
one-prefix upper counts for $w_*<t<h$. Here
\[
 T=t,\quad \nu=(1+t)/2,\quad L=(1-t)/2,\quad
 \chi=t,\quad s=L/t,\quad\alpha=2.
\]
The full square is $2t<\nu$, since $t<h<\mu<1/3$.
The selected largest-prime branch persists below $t$.
The cutoff equal to the marked prime is a permitted structural equality:
every later sieve prime is strictly smaller. Hence
\begin{equation}\label{eq:L2}
 L_2^*=8\alow\left(\frac{1+h}{2h}\right)
   -\frac4{249/50}\alow(249/50)
   +4\int_{w_*}^{h}\frac{\HH((1-t)/(2t))}{t(1-t)/2}\,dt.
\end{equation}
The table argument ranges from $(1-h)/(2h)>2$ to $199/50$, not
$249/50$; the latter belongs to the starting lower term.
For a cell $[q_i,q_{i+1}]$, its exact contributing interval is
\[
 (w_*,h)\cap[1/(2q_{i+1}+1),1/(2q_i+1)].
\]
The contribution to the Chen coefficient is $H_i$ times the integral
of $2/[t(1-t)]$, namely a difference of $2\log(t/(1-t))$.

For $U_3$ put $\nu(t)=v(t)$ on $(a,e)$ and $\nu(t)=1/2$ on
$(e,1/3)$; for $U_4$ use $v(t)$ on $(a,d)$. The enhanced part has
$2t<v(t)$, $L=v(t)-t$, $\chi=t$, and $\alpha\ge2$ because
$v(t)\le(1+t)/2$. Use the physical cutoff
\[
 k(t)=\min\{a,L/2\},\qquad s(t)=\max\{2,L/a\}.
\]
It is below the marked slot. The classical part is obtained from
$\nu<1/2$, with $\alpha>2$ and the same cutoff comparison.
Consequently
\begin{equation}\label{eq:U34}
 U_j^*=4\int_{I_j}\frac{A(s(t))-\HH(s(t))}{tL(t)}\,dt,
 \quad I_3=(a,1/3),\ I_4=(a,d).
\end{equation}
The upper defect is zero if $s>5$. Because $A=1$ up to three,
the cutoff lowering leaves the ordinary kernel unchanged.
On each affine piece $L=A_0-B_0t$, an unclamped cell has inverse
image $A_0-aq_{i+1}\le B_0t\le A_0-aq_i$; the clamped region is
$L\le2a$. The positive correction primitive is exactly
\[
 \int\frac{dt}{t(A_0-B_0t)}
 =\frac1{A_0}\log\frac{t}{A_0-B_0t}.
\]
For the original range, the support begins at $t\ge1-10a$ on the
first affine piece, intersected with that piece and the actual domain.
These cell cuts also supply a simple exact alternative to integrating
a discontinuous lookup without splitting.

\section*{The two-prefix lower terms}
Write the marked slots in decreasing order as $(t,u)$ and retain
\[
 V=\begin{cases}
 \max\{v(t),v(t+u),\Lam-3a/8\},&t<\mu,\\
 v(t),&t>\mu,
 \end{cases}\qquad L=V-t-u,\qquad s=L/a.
\]
The domains are $D_5=\{a<u<t<h\}$ and
$D_6=\{a<u<h<t<d\}$. The physical cutoff is always $a$.
Whenever $s\ge2$, the original full tests follow from
\[
 2t<v(t)\le V,\qquad t+2u<V,
 \quad\text{since }L\ge2a>h\ge u.
\]
The $v(t)$ branch keeps the largest slot and ceiling $u$;
the $v(t+u)$ branch keeps the leading selected pair and ceiling $u$.
The constant branch keeps its future third slot at most $a$ and
uses ceiling $a$. Its branch gate is $v(t)$, so use it above that
gate only when it actually improves the envelope; ties may use
another winning branch. All residual and newly appended primes lie
below the retained ceiling. Short residual atoms retain the original
classical fallback; the $W$ expression is not assigned an artificial
third slot when none exists.

For the first two branches $2V\le1+t+u$. For the constant branch
this holds automatically when $a\ge25/352$, since
\[
 \Lam-3a/8\le1/2+a\le(1+t+u)/2.
\]
Thus $\alpha\ge2$ on the original $a>1/12$ domain. Moreover
\[
 \frac{V-t-u}{a}\le\max\left\{
 \frac1{2a}-\frac32,\ \frac1{2a}-1,\ \frac\Lam a-\frac{19}8
 \right\}=\frac1{2a}-1<5.
\]
Every positive lower region is therefore within the certified table.
The legal replacements are
\begin{equation}\label{eq:L56}
 L_j^*=4\iint_{D_j}\frac{\alow(L/a)+\hh_{\rm adm}(t,u)}{tuL}
 \,du\,dt,\qquad j=5,6,
\end{equation}
where $\hh_{\rm adm}=\hh(L/a)$ only if the full state is admitted
and $2\le L/a\le5$, and is zero otherwise. Below two use
nonnegativity and no enhanced-state claim. For the positive-$a$
extension also check $2V\le1+t+u$ explicitly: on a winning constant
branch this is $t+u\ge\mu-3a/4$. Failure leaves the established
ordinary lower term unchanged. A comparison at a different level
would be a competing estimate, not an additive correction to this one.

On a region with a fixed affine winner $V$, cell membership is the
pair of linear inequalities $aq_i\le V-t-u\le aq_{i+1}$.
Together with ordering, full-square, ratio and branch gates, these
give an implementable rational polygon subdivision. Coarser boxes
may discard uncertain defects while retaining ordinary lower bounds.

\section*{The four-prefix option and the untouched switching terms}
For $j=10,11$ write $(x,y,z,w)$ in decreasing order, put $T=x+y+z+w$,
and use the full domains
\[
 D_{10}=\{a<w<z<y<x<h\},\quad
 D_{11}=\{a<w<z<y<h<x<d\},\quad d\mu=\frac{dx\,dy\,dz\,dw}{xyzw}.
\]
The original quotient cutoff is $z$. An ordinary comparison lowers
it to $\min(w,(\nu-T)/2)$; thus the specified four slots are largest
in its modulus atom. They need not be largest in the original integer.
Use exactly the gated cap set of \texttt{four-prefix-safe-minimum.tex}:
\[
 W=\max(2x,x+2y,x+y+2z,x+y+z+2w),
\]
and for each nonempty subset $I$ put
\[
 s_I=\sum_{i\in I}t_i,\qquad
 G_I=\max_{j\notin I}\left(s_I+
              \sum_{k<j,\ k\notin I}t_k+2t_j\right).
\]
The selected cap $v(s_I)$ requires $v(s_I)>\max(1/2,W,G_I)$.
The extra $B(z)=\Lam-3z/8$ cap requires
$x<\mu$ and $B(z)>\max(1/2,W,v(x))$.
Freeze a branch and a level strictly between its gate and cap on a
common positive box; only afterwards approach the endpoint.
Keep every selected slot and all four full tests, with ceiling $w$.
The old first three slots of a $B(z)$ branch persist below that ceiling.
Also $\alpha\ge2$, since $T>4a>1/3$ and $\nu\le\Lam$.
The positive-$a$ extension still implies $4a>25/112>2\Lam-1$,
so this ratio guard remains automatic there.

For an eligible endpoint $V_4$, put $L=V_4-T$,
$s=\max(2,L/w)$ and
\[
 O_0=4A(s)/L,\qquad O_H=4\{A(s)-\HH(s)\}/L,
 \qquad S_K=K\omega((1-T)/z)/z.
\]
Else put $O_0=O_H=+\infty$. The upper bound is
\begin{equation}\label{eq:four}
 U_j^*=\int_{D_j}\min(S_K,O_H)\,d\mu.
\end{equation}
Relative to an existing ordinary minimum its gain in the Chen
coefficient is exactly
$\frac14\int[\min(S_K,O_0)-\min(S_K,O_H)]d\mu$.
It is not the sum of switching and ordinary improvements.
The earlier proof gives $L/w<3$ for $a>1/12$, so $A(s)=1$ there;
keep the full $A$ in the positive-$a$ extension. Off the enhanced
eligible set retain switching (or a separately proved classical
ordinary option, without claiming it is supplied by this cap set).

No new defect is inserted into the switched $U_7,U_8,U_9$.
Their positive-product bounds retain their exact quotient domains;
in particular the moving square-root cutoff and all multiplicities
of $Y_9$ are unchanged. The complete functional is
\[
 c^*=\frac{3L_1+L_2^*+L_5^*+L_6^*-U_3^*-U_4^*
                    -2U_7-U_8-U_9-U_{10}^*-U_{11}^*}{4}.
\]
The table was built with $K_0=6.81$. A smaller separately proved
switching coefficient may be used for the five original switched
terms without recomputing or increasing these stored defects.
Old narrow-seed gains already present in a previous coefficient
must not also be added on overlapping domains: rebuild the ordinary
terms with these formulas, or take the maximum of separately proved
defects at exactly the same state.

\section*{The conservative one-dimensional switching option}
The new arithmetic script instead keeps switching alone for $U_{10}$
and $U_{11}$. Let $C=70893/125000$. If
\[
 C(1/h-4)>1,\qquad C((1-d)/h-3)>1,
\]
then the respective Buchstab arguments exceed $1/C$.
The proved tail bound $\omega\le C$ and positive integration give
\begin{align*}
 U_{10}&\le\frac{KC}{2}\int_a^h
       \frac{\log^2(h/t)\log(t/a)}{t^2}\,dt,\\
 U_{11}&\le KC\log(d/h)\int_a^h
       \frac{\log(h/t)\log(t/a)}{t^2}\,dt.
\end{align*}
Indeed fix $z=t$; integration in $w$ gives $\log(t/a)$.
For $U_{10}$ the ordered $z<y<x<h$ integral gives
$\log^2(h/z)/2$; for $U_{11}$ the independent $h<x<d$
integral gives $\log(d/h)$ and $z<y<h$ gives $\log(h/z)$.
The full fourth-prime interval is retained.

If additionally $d\ge1/5$, change variables in $U_7$ by
$r=(1-t-u)/u$. Then $1\le r\le R=1/t-2\le3$ and
\[
 \int_1^R\omega(r)\,dr=\log R+A(R+1)-1,
\]
by the explicit Buchstab formula up to three and the delay formula
for $A$ up to four. Hence the exact expression is
\[
 U_7=K\int_d^{1/3}
  \frac{\log((1-2t)/t)+A(1/t-1)-1}{t(1-t)}\,dt.
\]
For $d<1/5$ this reduction requires additional Buchstab pieces;
the displayed formula must not be extended automatically.
The unchanged formulas are
\[
 U_8=K\int_a^{1/3}\frac{\log(2-3t)}{t(1-t)}dt,\qquad
 U_9=K\int_h^d\frac{\log((1-t-d)/d)}{t(1-t)}dt.
\]

For $(a,h,d)=(417/5000,19/125,59/250)$ all stated gates hold.
The $L_2$ argument runs through $[53/19,199/50]$;
the first single-prefix supported point is $83/500$;
$U_4$ ends at $s=12025/3336$;
the classical $U_3$ clamp starts at $t=833/2500$.
The pair argument is at most $2083/417<5$ and the eligible
four-prefix argument is at most $1249/417<3$ before clamping.
The two uniform tail lower bounds are $49/19$ and $77/38$,
both larger than $1/C$.
For the older fixed tuple the latter coarse tail lower bound is
only $65247/38300<1/C$; these two convenient tail formulas are
therefore not an unconditional simplification for every threshold tuple.

\section*{Limits, implementation audit, and remaining scope}
First fix compact positive boxes, branch choices and buffered levels.
The typed-table theorem holds on each such averaged box, not for each
individual prime tuple. Refine a common partition for all choices,
take $N\to\infty$, and then exhaust strict domains while shrinking
buffers. Ordering, nontrivial affine gate and table-cell boundaries
have zero integral measure. A clamped argument identically two is
not discarded: the actual cutoff is $L/2$, and its closed-cell bound
applies by the endpoint convention established in the closure audit.
The same applies to a cutoff structurally equal to a marked prime.
All relevant residuals remain bounded below on fixed boxes and all
recursive depths are finite. No uniform assertion as the initial
residual tends to zero is needed.

No domain counterexample was found to the state selections in
\texttt{probe\_new\_finite\_double\_sieve.py} within its explicitly
asserted threshold range. Its pair ratio mask is necessary in the
positive-$a$ extension; its other pair guards follow from $h<2a$
where the lower defect is nonzero. Its four-prefix ratio condition
follows from $a>w_*/2$, even though not coded explicitly. Its selected
subset tests and the external full-square test together supply the
stated gates. This does not validate its interpolated values, larger
operator bank, quadrature, or optimization as a certificate.

For completeness, the probe's older three-slot envelope below $x=\mu$
contains apparently unfiltered pair caps. They cannot invalidate the
maximum: if $v(s)>v(x)$ and the selected pair has exponent $s$, its
additional gate is at most $2x+s$. For $s\le\mu$,
$v(s)-(2x+s)=(1-4x-s)/2\ge3/256$.
For $s>\mu$, the improvement condition gives
$s<25/96-x/3$, and hence
$v(s)-(2x+s)>23/96-7x/6\ge3/256$.
If the cap does not improve $v(x)$, use the already admitted largest
slot branch. The other prefix-product caps have no extra selected
gate; the $B(z)$ cap is used only when it exceeds $v(x)$.
Above $\mu$ the probe explicitly filters the surviving non-leading
branches by their selected gates. Its all-subset loop is followed by
the independent complete four-square check, as required here.

In particular ``look up only $s$'' is an invalid general rule:
$L_1$ gives an actual ratio failure. Also the retained-ceiling example
in the previous audit, with old slots $(.16,.15)$, $\nu=.581$ and
future-third bound $.04$, would admit an illegal new prime $.10569$
if that ceiling were forgotten. It is excluded by this prescription;
it is not a counterexample inside the probe's properly ordered domains.

For \texttt{certify\_finite\_table\_chen\_functional.py}, whole-box
minimum defects have the correct direction for a certified lower
coefficient after the signed terms are combined. Printed upper
endpoints using these minimum defects are not automatically upper
bounds for the exact pointwise step-table functional, nor for actual
prime counts. They must not support a claim of optimality or impossibility.
This note audits the default tuple and formulas, not that script's
large integration output or all interval-arithmetic implementation.
The remaining numerical task is an outward integration and independent
replay of this admitted functional; the primitive analytic distribution
dependencies remain separate from both local PASS results.
\end{document}
